\documentclass[12pt,a4paper]{article}
\usepackage[utf8]{inputenc}
\usepackage[spanish]{babel}
\usepackage{amsmath}
\usepackage{amsfonts}
\usepackage{amssymb}
\usepackage{graphicx}
\usepackage{tikz}
\usepackage{listings}
\usepackage{xcolor}
\usepackage[margin=2.5cm]{geometry}
\usepackage{hyperref}

\usetikzlibrary{arrows.meta, positioning, shapes}

\title{\textbf{Estructuras de Datos y Archivos}\\
\large Resumen Unidad 1: Modelos Matemáticos y Representaciones}
\author{Universidad Católica de Salta}
\date{\today}

\begin{document}

\maketitle
\tableofcontents
\newpage

\section{¿Qué es una Estructura de Datos?}

Una \textbf{estructura de datos} es simplemente una forma de organizar información en la computadora.

\subsection{Ejemplos Cotidianos}
\begin{itemize}
    \item \textbf{Agenda telefónica}: nombres y teléfonos organizados
    \item \textbf{Árbol genealógico}: muestra quién es hijo/padre de quién
    \item \textbf{Organigrama de empresa}: muestra quién es jefe de quién
\end{itemize}

\subsection{Ejemplo Visual: Organigrama}

\begin{center}
\begin{tikzpicture}[
    node distance=1.5cm and 2cm,
    every node/.style={rectangle, draw, rounded corners, minimum width=2.5cm, minimum height=1cm, align=center}
]
    \node (pres) {Presidente};
    \node (gerA) [below left=of pres] {Gerente A};
    \node (gerB) [below right=of pres] {Gerente B};
    \node (emp1) [below=of gerA] {Empleado 1};
    \node (emp2) [below=of gerB] {Empleado 2};
    
    \draw[-{Latex[length=3mm]}] (pres) -- (gerA);
    \draw[-{Latex[length=3mm]}] (pres) -- (gerB);
    \draw[-{Latex[length=3mm]}] (gerA) -- (emp1);
    \draw[-{Latex[length=3mm]}] (gerB) -- (emp2);
\end{tikzpicture}
\end{center}

\section{Grafos: El Modelo Matemático Básico}

Un \textbf{grafo} es como un mapa de conexiones. Tiene dos partes:

\begin{enumerate}
    \item \textbf{Puntos (P)} = ``Nodos'' = lugares/elementos
    \item \textbf{Flechas (E)} = ``Arcos'' = conexiones entre lugares
\end{enumerate}

\subsection{Ejemplo Visual de un Grafo}

\begin{center}
\begin{tikzpicture}[
    node distance=2.5cm,
    every node/.style={circle, draw, minimum size=1cm}
]
    \node (A) {A};
    \node (B) [right=of A] {B};
    \node (C) [below=of A] {C};
    \node (D) [right=of C] {D};
    
    \draw[-{Latex[length=3mm]}] (A) -- (B);
    \draw[-{Latex[length=3mm]}] (A) -- (C);
    \draw[-{Latex[length=3mm]}] (B) -- (C);
    \draw[-{Latex[length=3mm]}] (C) -- (D);
\end{tikzpicture}
\end{center}

\textbf{Notación matemática:}
\begin{align*}
G &= (P, E) \\
P &= \{A, B, C, D\} \\
E &= \{(A,B), (A,C), (B,C), (C,D)\}
\end{align*}

\begin{itemize}
    \item Desde A puedes ir a B
    \item Desde A puedes ir a C
    \item Desde B puedes ir a C
    \item Desde C puedes ir a D
\end{itemize}

\section{Conceptos Importantes}

\subsection{Minimal (Nodo de Entrada)}

Un nodo es \textbf{minimal} cuando nadie apunta hacia él (excepto posiblemente él mismo).

\begin{center}
\begin{tikzpicture}[
    node distance=2cm,
    every node/.style={circle, draw, minimum size=1cm}
]
    \node (A) {A};
    \node (B) [right=of A] {B};
    \node (C) [right=of B] {C};
    
    \draw[-{Latex[length=3mm]}] (A) -- (B);
    \draw[-{Latex[length=3mm]}] (B) -- (C);
    
    \node[below=0.3cm of A, draw=none] {\textcolor{red}{MINIMAL}};
\end{tikzpicture}
\end{center}

\subsection{Maximal (Nodo Final)}

Un nodo es \textbf{maximal} cuando él no apunta a nadie (excepto posiblemente a sí mismo).

\begin{center}
\begin{tikzpicture}[
    node distance=2cm,
    every node/.style={circle, draw, minimum size=1cm}
]
    \node (A) {A};
    \node (B) [right=of A] {B};
    \node (C) [right=of B] {C};
    
    \draw[-{Latex[length=3mm]}] (A) -- (B);
    \draw[-{Latex[length=3mm]}] (B) -- (C);
    
    \node[below=0.3cm of C, draw=none] {\textcolor{red}{MAXIMAL}};
\end{tikzpicture}
\end{center}

\newpage
\subsection{Lazo}

Una flecha que vuelve al mismo punto.

\begin{center}
\begin{tikzpicture}[
    node distance=2cm,
    every node/.style={circle, draw, minimum size=1cm}
]
    \node (A) {A};
    \node (B) [right=of A] {B};
    \node (C) [below=of B] {C};
    
    \draw[-{Latex[length=3mm]}] (A) -- (B);
    \draw[-{Latex[length=3mm]}] (B) -- (C);
    \draw[-{Latex[length=3mm]}] (A) to [out=135, in=225, looseness=5] (A);
    
    \node[left=0.3cm of A, draw=none] {\textcolor{red}{LAZO}};
\end{tikzpicture}
\end{center}

\subsection{Right(x) - ¿A quiénes apunto desde X?}

\begin{center}
\begin{tikzpicture}[
    node distance=2cm,
    every node/.style={circle, draw, minimum size=1cm}
]
    \node (A) {A};
    \node (B) [above right=of A] {B};
    \node (C) [below right=of A] {C};
    
    \draw[-{Latex[length=3mm]}] (A) -- (B);
    \draw[-{Latex[length=3mm]}] (A) -- (C);
\end{tikzpicture}
\end{center}

$$\text{Right}(A) = R(A) = \{B, C\}$$

\subsection{Left(x) - ¿Quiénes apuntan hacia X?}

\begin{center}
\begin{tikzpicture}[
    node distance=2cm,
    every node/.style={circle, draw, minimum size=1cm}
]
    \node (C) {C};
    \node (A) [above left=of C] {A};
    \node (B) [below left=of C] {B};
    
    \draw[-{Latex[length=3mm]}] (A) -- (C);
    \draw[-{Latex[length=3mm]}] (B) -- (C);
\end{tikzpicture}
\end{center}

$$\text{Left}(C) = L(C) = \{A, B\}$$

\subsection{Input Degree - Flechas que LLEGAN}

\begin{center}
\begin{tikzpicture}[
    node distance=2cm,
    every node/.style={circle, draw, minimum size=1cm}
]
    \node (C) {C};
    \node (A) [above left=of C] {A};
    \node (B) [below left=of C] {B};
    
    \draw[-{Latex[length=3mm]}] (A) -- (C);
    \draw[-{Latex[length=3mm]}] (B) -- (C);
\end{tikzpicture}
\end{center}

$$\text{Input degree}(C) = |L(C)| = 2$$

\subsection{Output Degree - Flechas que SALEN}

\begin{center}
\begin{tikzpicture}[
    node distance=2cm,
    every node/.style={circle, draw, minimum size=1cm}
]
    \node (A) {A};
    \node (B) [above right=of A] {B};
    \node (C) [below right=of A] {C};
    
    \draw[-{Latex[length=3mm]}] (A) -- (B);
    \draw[-{Latex[length=3mm]}] (A) -- (C);
\end{tikzpicture}
\end{center}

$$\text{Output degree}(A) = |R(A)| = 2$$

\newpage
\section{Conceptos de Camino/Paso}

\subsection{Paso - Un camino de un nodo a otro}

\begin{center}
\begin{tikzpicture}[
    node distance=2cm,
    every node/.style={circle, draw, minimum size=1cm}
]
    \node (A) {A};
    \node (B) [right=of A] {B};
    \node (C) [right=of B] {C};
    \node (D) [right=of C] {D};
    
    \draw[-{Latex[length=3mm]}, thick, red] (A) -- (B);
    \draw[-{Latex[length=3mm]}, thick, red] (B) -- (C);
    \draw[-{Latex[length=3mm]}, thick, red] (C) -- (D);
\end{tikzpicture}
\end{center}

\textbf{Paso de A a D:} $A \rightarrow B \rightarrow C \rightarrow D$

\textbf{Longitud del paso:} 3 (tres flechas)

\subsection{Ciclo - Camino que vuelve al inicio}

\begin{center}
\begin{tikzpicture}[
    node distance=2.5cm,
    every node/.style={circle, draw, minimum size=1cm}
]
    \node (A) {A};
    \node (B) [right=of A] {B};
    \node (C) [below right=of A] {C};
    
    \draw[-{Latex[length=3mm]}, thick, blue] (A) -- (B);
    \draw[-{Latex[length=3mm]}, thick, blue] (B) -- (C);
    \draw[-{Latex[length=3mm]}, thick, blue] (C) -- (A);
\end{tikzpicture}
\end{center}

\textbf{Ciclo:} $A \rightarrow B \rightarrow C \rightarrow A$

\section{Representaciones en la Computadora}

\subsection{Representación Secuencial (con arrays)}

\textbf{Matriz de Adyacencia} para el grafo: $A \rightarrow B \rightarrow C$

$$
\begin{array}{c|ccc}
  & A & B & C \\
\hline
A & 0 & 1 & 0 \\
B & 0 & 0 & 1 \\
C & 0 & 0 & 0 \\
\end{array}
$$

\begin{itemize}
    \item 1 = hay flecha
    \item 0 = no hay flecha
\end{itemize}

\subsection{Representación Linkeada (con punteros)}

Cada nodo guarda la dirección del siguiente:

\begin{center}
\begin{tikzpicture}[
    node distance=0.5cm,
    box/.style={rectangle, draw, minimum width=1.5cm, minimum height=0.8cm}
]
    \node[box] (A1) {A};
    \node[box, right=of A1] (A2) {$\bullet$};
    \node[box, right=1.5cm of A2] (B1) {B};
    \node[box, right=of B1] (B2) {$\bullet$};
    \node[box, right=1.5cm of B2] (C1) {C};
    \node[box, right=of C1] (C2) {NULL};
    
    \draw[-{Latex[length=3mm]}] (A2) -- (B1);
    \draw[-{Latex[length=3mm]}] (B2) -- (C1);
    
    \node[below=0.1cm of A1] {\tiny dato};
    \node[below=0.1cm of A2] {\tiny ptr};
    \node[below=0.1cm of B1] {\tiny dato};
    \node[below=0.1cm of B2] {\tiny ptr};
    \node[below=0.1cm of C1] {\tiny dato};
    \node[below=0.1cm of C2] {\tiny ptr};
\end{tikzpicture}
\end{center}

\begin{itemize}
    \item $\bullet$ = puntero (dirección de memoria)
    \item NULL = no apunta a nada
\end{itemize}

\section{Ejemplo Práctico: Red de Ciudades}

\begin{center}
\begin{tikzpicture}[
    node distance=3cm,
    every node/.style={rectangle, draw, rounded corners, minimum width=2cm, minimum height=1cm}
]
    \node (S) {Salta};
    \node (T) [right=of S] {Tucumán};
    \node (St) [below right=of S] {Santiago};
    
    \draw[-{Latex[length=3mm]}] (S) -- node[above] {300km} (T);
    \draw[-{Latex[length=3mm]}] (S) -- node[left] {400km} (St);
    \draw[-{Latex[length=3mm]}] (T) -- node[right] {200km} (St);
\end{tikzpicture}
\end{center}

\begin{align*}
P &= \{\text{Salta}, \text{Tucumán}, \text{Santiago}\} \\
E &= \{(\text{Salta}, \text{Tucumán}, 300), \\
  &\quad (\text{Salta}, \text{Santiago}, 400), \\
  &\quad (\text{Tucumán}, \text{Santiago}, 200)\}
\end{align*}

\textbf{Análisis:}
\begin{itemize}
    \item Right(Salta) = \{Tucumán, Santiago\}
    \item Left(Santiago) = \{Salta, Tucumán\}
    \item Paso directo Salta $\rightarrow$ Santiago = 400km
    \item Paso Salta $\rightarrow$ Tucumán $\rightarrow$ Santiago = 500km
\end{itemize}

\end{document}